% Successor-side research note based on analysis/tumbling/study.md.
\subsection{Ballistic Tumbling and Detached-Body Dynamics}\label{sec:ballistic-tumbling}

This subsection records the successor-side mathematics used for detached
aero-ballistic bodies. It is an engineering model and research note, not a
claim about an undocumented historical TAOS equation. The working source is
the project note \texttt{analysis/tumbling/study.md}; calibrated coefficient
data must replace the illustrative closures before a vehicle-specific claim is
made.

For a body moving through air, dynamic pressure and drag area provide a useful
separation between the environment and the aerodynamic closure:
\[
  q_\infty=\frac{1}{2}\rho\lVert\vec V_{w}\rVert^2,
  \qquad
  D=q_\infty C_D S_{\rm ref}=q_\infty K_D,
  \qquad
  K_D=C_D S_{\rm ref}.
\]
The drag-area form is important when comparing sources that use different
reference areas. Projected area is a first-order geometry model; the actual
coefficient may also depend on Mach number, Reynolds number, base pressure,
surface condition, and the orientation of the relative wind.

Let $\hat x_b$ be the longitudinal body axis and let $\alpha$ be the angle
between $\hat x_b$ and the relative-wind direction. A useful cylinder
projected-area closure is
\[
  S_{\rm proj}(\alpha)
  =\pi r^2\lvert\cos\alpha\rvert
   +2rL\lvert\sin\alpha\rvert,
\]
where $r$ is radius and $L$ is length. A spheroid uses its axial and
transverse semi-axes in the corresponding orientation-dependent ellipse
closure. Cones require separate nose-first, broadside, and base-first anchors;
interpolation between those anchors must not be presented as measured data.

For a freely tumbling rigid body, the body-rate equation is
\[
  \mathbf I\dot{\boldsymbol\omega}
  +\boldsymbol\omega\times
    (\mathbf I\boldsymbol\omega)
  =\mathbf M_{\rm aero}+\mathbf M_{\rm sep}+\mathbf M_{\rm control}.
\]
If the aerodynamic force is applied at a center-of-pressure offset, its
moment contribution is
\[
  \mathbf M_{\rm aero}
  =(\mathbf r_{\rm cp}-\mathbf r_{\rm cg})
    \times\mathbf F_{\rm aero}
  +\mathbf M_{\rm damping}.
\]
The Alpha 3 ballistic surrogate makes the unresolved rate damping explicit.
That damping stabilizes coarse envelope integrations; it is not a substitute
for a measured or computed aerodynamic moment database.

Quaternion attitude propagation uses the body rate without Euler-angle
singularities:
\[
  \dot q=\frac{1}{2}q\otimes
  \begin{bmatrix}0\\\boldsymbol\omega\end{bmatrix},
  \qquad \lVert q\rVert=1.
\]
The implementation renormalizes the integrated quaternion and records the
attitude, body rates, projected area, force, moment, and termination state in
the child telemetry.

For isotropic tumbling studies, an orientation-average coefficient is defined
over the unit sphere by
\[
  \overline{K_D}
  =\frac{1}{4\pi}\int_{S^2}K_D(\hat V_w)\,d\Omega.
\]
This average is useful for comparison and screening, but it must not replace
the time history when attitude, impact timing, or event coupling matters.
The computational fidelity ladder therefore remains explicit: point-mass
3-DOF, pseudo-6-DOF attitude/tumble surrogates, and native rigid-body 6-DOF.
